% Diplomatic French for printed p.221.
% Nine print-level defects are preserved for the corrected layer to repair.
\clearpage
\SourcePage{221}

\LayerTarget{tohoku:bibliography}
\begin{center}\normalsize\textsc{Bibliographie}\end{center}
\addcontentsline{toc}{chapter}{Bibliographie}
\markboth{Bibliographie}{Bibliographie}

\small
\begin{list}{}{%
  \setlength{\leftmargin}{5.3em}%
  \setlength{\itemindent}{-3em}%
  \setlength{\labelwidth}{1.6em}%
  \setlength{\labelsep}{0.7em}%
  \setlength{\itemsep}{0.28em}%
  \setlength{\parsep}{0pt}%
  \setlength{\topsep}{0.3em}%
}
\LayerTarget{tohoku:bib:1}
\item[{[1]}] M.\ F.\ \textsc{Atiyah}, Complex analytic connections,
Trans. Amer. Mat\SourceTurnedH{}. Soc.85(1957) 181--207.

\LayerTarget{tohoku:bib:2}
\item[{[2]}] A.\ \textsc{Borel}, Nouvelle démonstration d'un théorème de
P.\ A. Smith, Commentarii Math. Helv.29 (1955), 27--39

\LayerTarget{tohoku:bib:3}
\item[{[3]}] D.\ A.\ \textsc{Buchsbaum}, Exact categories and duality,
Trans. Amer. Math. Soc.80 (1955), 1--34

\LayerTarget{tohoku:bib:4}
\item[{[4]}] H.\ \textsc{Cartan}, Séminaire E.\ N.\ S. 1950/51.

\LayerTarget{tohoku:bib:5}
\item[{[5]}] H.\ \textsc{Cartan}, Séminaire E.\ N.\ S. 1951/52.

\LayerTarget{tohoku:bib:5-bis}
\item[{[5]}] bisH.\ \textsc{Cartan-C.\ Chevalley}, Séminbire E.\ N.\ S. 1955/56.

\LayerTarget{tohoku:bib:6}
\item[{[6]}] H.\ \textsc{Cartan-S.\ Eilenberg}, Homological Algebra,
Princeton University Press 1956.

\LayerTarget{tohoku:bib:7}
\item[{[7]}] H.\ \textsc{Cartan-J.\ P. Serre}, Un théorème de finitude,
C.R. Acad. Sci. Paris 237 (1953), 128--130.

\LayerTarget{tohoku:bib:8}
\item[{[8]}] P.\ \textsc{Cartier} (Référence non parvenue)

\LayerTarget{tohoku:bib:9}
\item[{[9]}] R.\ \textsc{Godement}, Théorie des faisceaux, en préparation.

\LayerTarget{tohoku:bib:10}
\item[{[10]}] R.\ \textsc{Godement}, Cohomologie des groupes discontinus,
Séminaire Bourbaki, Mars 1954.

\LayerTarget{tohoku:bib:11}
\item[{[11]}] A.\ \textsc{Grothendieck}, A general theory of fibre spaces
with structure sheaf, University of Kansas 1955.

\LayerTarget{tohoku:bib:12}
\item[{[12]}] G.\ \textsc{Hochschild}, Group extensions of Lie groups I,
Ann. of Math, 54(1951), 96--109

\LayerTarget{tohoku:bib:13}
\item[{[13]}] G.\ \textsc{Hochschild}, Group extensions of Lie groups II,
Ann. of Math, 54--1951), 537--551.

\LayerTarget{tohoku:bib:14}
\item[{[14]}] S.\ \textsc{MacLane}, Duality for groups,
Bull. Amer.Math. Soc., 56(1950), 485--516.

\LayerTarget{tohoku:bib:15}
\item[{[15]}] J.\ P.\ \textsc{Serre}, Faisceaux algébriques cohérents,
Ann. of Math., 61 (1955),197--278.

\LayerTarget{tohoku:bib:16}
\item[{[16]}] J.\ P.\ \textsc{Serre}, Géométrie Analytique et Géométrie
Algébrique. Anna Inst.Fourier VI(1955--56)1--42.

\LayerTarget{tohoku:bib:17}
\item[{[17]}] J.\ P.\ \textsc{Serre}, Classes de groupes abéliens et groupes
d'homotopie, Ann. of Math., 58(1953), 258--294.

\LayerTarget{tohoku:bib:18}
\item[{[18]}] J.\ P.\ \textsc{Serre}, Sur la cohomologie des variété
algébriques, Journ. Math. pures et appl. 36(1957)\mbox{1--16}.

\LayerTarget{tohoku:bib:19}
\item[{[19]}] A.\ \textsc{Shapiro}, Ann. of Math., 50 (1949), 581--586.

\LayerTarget{tohoku:bib:20}
\item[{[20]}] A.\ \textsc{Weil}, Fibre spaces in Algebraic Geometry
(Notes by A. Wallace), Chicago 1952.

\LayerTarget{tohoku:bib:21}
\item[{[21]}] A.\ \textsc{Weil}, Généralisation des fonctions abéliennes,
J.Math. Pures Appl. 17(1938), 47--87.
\end{list}
\normalsize

\TohokuTerminal{221}
